\section{Credenziali di accesso}
Sul sito web è possibile accedere come utente normale o amministratore:
\begin{center}
\begin{tabular}{|l|c|c|}
\hline
\rowcolor{brown}

\color{white}\textbf{Area} &
\color{white}\textbf{Username} &
\color{white}\textbf{Password} \\
\rowcolor{beige} 
\hline
\textbf{Area utente} & \texttt{user} & \texttt{user} \\
\rowcolor{beige} 
\hline
\textbf{Area amministratore} & \texttt{admin} & \texttt{admin} \\
\hline
\end{tabular}
\end{center}
Si specifica che all'utente \texttt{user} è stata limitata la funzionalità di modifica password e di eliminazione account, per garantire che tutti possano sempre accedere tramite queste credenziali (fare la registrazione se si desidera provare queste funzionalità).\\
Inoltre si fa notare che il sito consente più admin, suggeriamo quindi un secondo account di test, per dimostrare questa funzionalità (poichè gli amministratori, essendo account "di lavoro", non vengono inseriti tramite registrazione).
\begin{center}
\begin{tabular}{|l|c|c|}
\hline
\rowcolor{brown}
\color{white}\textbf{Area} &
\color{white}\textbf{Username} &
\color{white}\textbf{Password} \\
\rowcolor{beige} 
\hline
\textbf{Area amministratore} & \texttt{ombretta.gaggi.admin@unipd.it} & \texttt{Gaggi!00} \\
\hline
\end{tabular}
\end{center}
\begin{mdframed}[
    backgroundcolor=gray!15,
    linecolor=none,
    leftmargin=-3cm,
    rightmargin=-3cm,
    innerleftmargin=4cm,
    innerrightmargin=4cm,
    innertopmargin=10pt,
    innerbottommargin=10pt
    ]
\textmd{Nota: una richiesta di adozione viene valutata dallo specifico amministratore che ha in carico l'animale. Si consiglia quindi di controllare gli animali assegnati agli account admin dati (tramite area amministratore) per procedere con una richiesta di adozione.}
\end{mdframed}
